\section{Limitations and next steps}

Ctl-2 studies an additive edit to a final residual with affine readouts. It does
not rerun the Transformer, and its nuisance measure is a fitted probe. The four
binary inputs yield four trajectories per trained model, so training seed is
the statistical unit. Affine-span projection is only a geometry diagnostic.
These choices make the gain and bias mechanisms inspectable, but they limit
claims about semantic behavior under real token-level interventions.

The IOI analysis uses a known small circuit and one intervention convention.
The current pair ranking is confounded by unequal group sizes and unequal
functional form. The next analysis must match pair-model capacity, report
leave-one-pair-out contributions, attach prompt-level uncertainty to direct
M\"obius effects, and quantify the power of the random-subset design for the
known whole-group interaction.

Artifact adoption also remains bounded. A versioned contract and one external
adapter are a necessary first step, not evidence that researchers will adopt a
general benchmark. The stronger adoption gate is a second adapter on a trained
or IOI task together with baseline observers from the literature. Those are
future extensions rather than claims of the v0 release.

The cyber fixture should not be iterated by rewriting the same intent clauses.
A useful future security task would make authorization depend on context, use
policy or rule systems as first-class baselines, and break lexical shortcuts
through cross-context splits.

\section{Conclusion}

ObserverBench asks a systems question: which estimate and actuation direction
is a controller allowed to trust? The answer is not global. In first-order
regimes a cheap estimator may suffice. In one-shot interactional regimes its
direction can move collateral state. In an affine closed loop, estimation bias
can drive sustained target-error growth even when direction is held fixed.
Response calibration can repair the dynamic gain and still leave the initial
state wrong.

The workbench must therefore report more than a method name. It must say what
was estimated, which factors changed, what the intervention distribution
covered, how well the measurements conditioned the estimate, and whether the
estimate was adequate for the stated use. The Phase 0--1 results establish that
discipline for the control tasks and make the artifact boundary executable. The
remaining IOI analysis asks whether the same discipline survives a fair
comparison of interaction models.
